\documentclass[11pt]{article}
\usepackage{amsmath,amssymb,amsthm}
\usepackage[margin=1.1in]{geometry}
\newtheorem{theorem}{Theorem}
\newtheorem{definition}{Definition}
\newtheorem{lemma}{Lemma}
\newtheorem{proposition}{Proposition}
\newtheorem{corollary}{Corollary}
\theoremstyle{remark}
\newtheorem{remark}{Remark}
\newtheorem{question}{Question}
\newcommand{\N}{\mathbb{N}}
\newcommand{\Q}{\mathbb{Q}}
\newcommand{\Z}{\mathbb{Z}}
\newcommand{\m}{\mathfrak{m}}
\newcommand{\A}{\mathcal{A}}

\title{The multiplication crystal's engine is a formal logarithm}
\author{Carl Gribble\thanks{Programme Note 2, draft v0.1. Companion to
\cite{CompanionI,CompanionII,CompanionIII}, the essay \cite{Essay}, and Programme Note 1
\cite{Defs}, whose vocabulary is used throughout. Developed in collaboration with Claude
(Anthropic); the machine checks of Section~7 were run as described there. Responsibility for the
content rests with the author. A full disclosure of AI use appears at the end of the note.}}
\date{August 2026 --- draft}

\begin{document}
\maketitle

\begin{abstract}
The identity engine of \cite{CompanionI} --- Linnik's alternating-harmonic de-duplication of the
multiplication table --- is restated and proved as a statement of pure algebra: in the convolution
algebra of any crystal (measured free monoid \cite{Defs}), the formal logarithm of the zeta element
is supported on axis powers with value $1/j$ at $p^{\,j}$. No limit, no analytic continuation, and
no property of the integers enters; the proof is unique factorization plus the exp/log calculus of
a complete filtered $\Q$-algebra, and is stated at Beurling generality. The M\"obius twin
$\Lambda = \ell * \mu$ is derived from the same calculus in two lines, and the classical toggle
involution gives it a fully bijective proof. For the harmonic form itself we record an honest
obstruction --- $\kappa$ is not integer-valued, so no unweighted sign-reversing involution can
compute it --- together with a Burnside reformulation in which the $1/k$ weights become orbit-size
corrections under cyclic rotation, verified in micro-examples and left open as the \emph{necklace
question}. Exact machine checks over random crystals (a basis-vector trick makes one run cover all
Beurling length systems at once) and over the integers confirm every identity at truncation. The
note closes with the revision plan this induces for \cite{CompanionI} and the statement that work
package WP7.H3 will mechanize.
\end{abstract}

\section{What is being claimed}

The engine of the companion paper \cite{CompanionI} is the per-integer identity
\[
\sum_{k \ge 1} \frac{(-1)^{k-1}}{k}\, D_k^*(n) \;=\; \frac{\Lambda(n)}{\log n},
\]
Linnik's identity \cite{Linnik}, there derived from the generating identity
$\log \zeta = \log(1 + (\zeta - 1))$ read coefficientwise. That derivation is correct but wears
analytic clothing: $\zeta(s)$ appears as a function of a complex variable, and the reader may
suspect that convergence, or some property of the primes, is doing work. The point of this note is
that nothing of the sort is involved. The identity is a theorem about \emph{free commutative
monoids with finite divisor sets}, provable by elementary algebra, valid verbatim for every crystal
of \cite{Defs} --- the integers, the function fields, every Beurling system --- and its natural
proof is three lines once the right algebra is named. Naming it also locates the identity's
\emph{bijective} content precisely: it lives in the M\"obius twin, while the harmonic weights
$1/k$ are calculus, not counting --- with one honest open door recorded in Section~6.

Throughout, $C = (M, \ell)$ is a crystal with axis set $P$ \cite{Defs}; the letter $p$ always
denotes an axis. The results of Sections 2--5 never use $\ell$ except where it is named
explicitly, and never use axiom (F); they are theorems about the bare free monoid $M$.

\section{The convolution algebra}

\begin{definition}
$\A(M)$ is the set of all functions $f \colon M \to \Q$ with pointwise addition and the convolution
\[
(f * g)(m) \;=\; \sum_{ab = m} f(a)\, g(b),
\]
the sum over ordered pairs with $ab = m$.
\end{definition}

The sum is finite: divisors of $m$ correspond to componentwise-smaller exponent vectors, of which
there are finitely many. $\A(M)$ is a commutative associative $\Q$-algebra with unit
$\delta = 1_{\{1\}}$. For $M = (\N_{\ge 1}, \times)$ it is the algebra of formal Dirichlet series;
in Rota's theory it is the reduced incidence algebra of the divisor lattice, translation-invariance
of intervals ($[a, b] \cong [1, b/a]$) being what ``reduced'' means \cite{Rota, DRS, StanleyEC}.

Write $\Omega \colon M \to \N$ for the dimension grading ($\Omega(p) = 1$, additive) and
\[
F_k \;=\; \{\, f \in \A(M) : f(m) = 0 \text{ whenever } \Omega(m) < k \,\}, \qquad
\m \;=\; F_1 = \{ f : f(1) = 0 \}.
\]

\begin{lemma}[Filtration]\label{lem:filt}
$F_j * F_k \subseteq F_{j+k}$; in particular $\m^{*k} \subseteq F_k$. Any series
$\sum_{k \ge 0} f_k$ with $f_k \in F_k$ defines an element of $\A(M)$ (at each $m$ only the terms
$k \le \Omega(m)$ contribute), and $\A(M)$ is complete for the topology this filtration defines.
\end{lemma}

\begin{proof}
If $ab = m$ with $f(a) g(b) \ne 0$ then $\Omega(m) = \Omega(a) + \Omega(b) \ge j + k$. The rest is
the definition of pointwise finiteness.
\end{proof}

Consequently the formal series
\[
\exp(f) = \sum_{k \ge 0} \frac{f^{*k}}{k!} \quad (f \in \m), \qquad
\log(\delta + f) = \sum_{k \ge 1} \frac{(-1)^{k-1}}{k}\, f^{*k} \quad (f \in \m)
\]
are well defined, and all the usual identities hold --- $\exp$ and $\log$ are mutually inverse
bijections between $\m$ and $\delta + \m$, and $\exp(f + g) = \exp f * \exp g$, equivalently
$\log(u * v) = \log u + \log v$ for $u, v \in \delta + \m$. (Proof in one sentence: in each
truncation $\A(M)/F_{n+1}$ the ideal generated by $\m$ is nilpotent, the identities are the
classical polynomial identities of exp/log there, and completeness assembles them; see
\cite[Ch.~5]{StanleyEC}.) Likewise, locally finite families multiply: if $(u_p)_{p \in P}$ satisfy
$u_p \in \delta + \m$ with $u_p - \delta$ supported on powers of $p$, then $\prod_p u_p$ is
defined (evaluation at $m$ stabilizes once the finite set of axes dividing $m$ is included) and
$\log \prod_p u_p = \sum_p \log u_p$.

Two more standing elements: $\zeta = $ the constant function $1$; $F = \zeta - \delta$, the
indicator of $M \setminus \{1\}$; and the layer counts $D_k := F^{*k}$, so that $D_k(m)$ is the
number of ordered $k$-tuples of non-identity elements with product $m$ --- the $k$-dimensional
table of \cite{CompanionI} at $C_{\Z}$.

\section{The formal logarithm theorem}

\begin{theorem}[The engine, at crystal generality]\label{thm:engine}
In $\A(M)$ define $\kappa = \log \zeta = \sum_{k \ge 1} \frac{(-1)^{k-1}}{k} D_k$. Then
\[
\kappa(m) \;=\;
\begin{cases}
1/j, & m = p^{\,j},\ p \in P,\ j \ge 1,\\
0, & \text{otherwise.}
\end{cases}
\]
\end{theorem}

\begin{proof}
For an axis $p$ let $e_p = 1_{\{p\}}$. Freeness gives $e_p^{*j} = 1_{\{p^{\,j}\}}$ (the only
ordered $j$-tuple of copies of $p$ with product $p^{\,j}$ is $(p, \dots, p)$), hence
\[
\sum_{j \ge 0} e_p^{*j} \;=\; 1_{\{p^{\,j} \,:\, j \ge 0\}} \;=\; (\delta - e_p)^{-1}.
\]
Unique factorization says exactly that $\zeta = \prod_{p} (\delta - e_p)^{-1}$: the coefficient of
$m$ in the finite sub-product over any axis set containing the support of $m$ is $1$. Taking
$\log$ of the locally finite product,
\[
\kappa \;=\; \sum_p \log\,(\delta - e_p)^{-1} \;=\; \sum_p \sum_{j \ge 1} \frac{e_p^{*j}}{j}
\;=\; \sum_p \sum_{j \ge 1} \frac{1}{j}\, 1_{\{p^{\,j}\}} . \qedhere
\]
\end{proof}

\begin{corollary}[Linnik's identity]\label{cor:linnik}
At $C_{\Z}$, Theorem~\ref{thm:engine} is Lemma~1 of \cite{CompanionI}:
$\sum_k \frac{(-1)^{k-1}}{k} D^*_k(n) = \Lambda(n)/\log n$. At $C_{\mathbb{F}_q[t]}$ it is the
function-field analogue, with no change in proof. Proposition~4 of \cite{Defs} is proved.
\end{corollary}

\begin{corollary}[Window pairings]\label{cor:pairing}
For any finitely supported weight $w \colon M \to \Q$,
\[
\sum_{m \in M} w(m)\, \kappa(m) \;=\; \sum_{p} \sum_{j \ge 1} \frac{w(p^{\,j})}{j}.
\]
At $C_\Z$ with $w(m) = |m|^{\,j}\,1_{\{|m| \le y\}}$ this is the identity
$\sum_{k} \frac{(-1)^{k-1}}{k} T_k^{(j)}(y) = \sum_{p^i \le y} p^{\,ij}/i$ that feeds the
extraction line of \cite{CompanionI}; the remaining content of that paper's formula system --- the
Faulhaber base, the hyperbola step, and the root-scale recursion --- is finite arithmetic
reorganization of these pairings, with no analytic input. The exactness assertion of
\cite{CompanionI} (``every harmonic combination lands on denominator $1$'') is the integrality of
the left side, now visible formally: it equals a finite sum of integers $p^{\,ij}$ divided by
$i$, summed against the M\"obius-inverted extraction, and the paper's recursion is precisely the
bookkeeping that clears the denominators.
\end{corollary}

\begin{corollary}[Duplicated-mass skeleton]\label{cor:dup}
$F - \kappa = \sum_{k \ge 2} \frac{(-1)^{k}}{k} D_k$ identically in $\A(M)$; evaluated at $m$,
this is $1 - 1/j$ on axis powers $p^{\,j}$ and $1$ on all other $m \ne 1$ --- the composite
indicator with prime-power dust of \cite[Thm.~2]{CompanionI}. The cancellation of the
two-dimensional table out of its own duplication ledger is the visible absence of $D_2$ from the
pair $(F, \kappa)$ once the $k = 2$ term is moved across; the rest of that theorem is the window
pairing of Corollary~\ref{cor:pairing} at weight $w(m) = |m|\,1_{\{|m| \le N\}}$ plus the fold
bookkeeping done there.
\end{corollary}

\section{The M\"obius twin, bijectively}

\begin{proposition}[M\"obius, at crystal generality]\label{prop:mu}
$\mu := \zeta^{-1} = \prod_p (\delta - e_p)$ is given by $\mu(m) = (-1)^{\omega(m)}$ if $m$ is a
product of distinct axes ($\omega$ of them) and $0$ otherwise; and $\sum_{d \mid m} \mu(d) = 0$
for every $m \ne 1$.
\end{proposition}

\begin{proof}
The product formula is the finite sub-product evaluation as in Theorem~\ref{thm:engine}; expanding
it gives the stated values. The divisor sum is $(\mu * \zeta)(m) = \delta(m)$. For the bijective
proof of that vanishing: fix an axis $p \mid m$; on the squarefree divisors of $m$ (the support of
$\mu$ among divisors), toggling the factor $p$ is an involution that reverses the sign of $\mu$
and has no fixed points. Mass cancels in pairs.
\end{proof}

Now bring in the length. Any additive weight defines a derivation: for
$L(f)(m) := \ell(m) f(m)$,
\[
L(f * g) \;=\; L(f) * g + f * L(g),
\]
because $\ell(m) = \ell(a) + \ell(b)$ on each pair $ab = m$; $L$ respects the filtration. The
standard consequence $L(\log u) = L(u) * u^{-1}$ (proved in the nilpotent truncations, as before)
yields, in two lines, the twin form of the engine:

\begin{proposition}[$\Lambda = \ell * \mu$]\label{prop:twin}
Let $\Lambda := L(\kappa)$, so $\Lambda(p^{\,j}) = \ell(p)$ on axis powers and $0$ elsewhere.
Then
\[
\Lambda \;=\; L(\zeta) * \zeta^{-1} \;=\; \ell * \mu ,
\]
where $\ell$ denotes the function $m \mapsto \ell(m)$. Equivalently $\ell = \Lambda * \zeta$,
which is the direct count $\sum_{p^{j} \mid m} \ell(p) = \ell(m)$.
\end{proposition}

The toggle involution of Proposition~\ref{prop:mu} extends to $\ell * \mu$ by linearity in the
$\ell$-leg, so the twin identity --- hence the engine's support statement, hence the composite
indicator used by the duplicated-mass theorem --- has a fully bijective proof. This is the
integer-weight ``dimension parity'' route that \cite{CompanionI} calls Legendre's twin; the
calculus above is the precise bridge between the two.

\section{Where the bijection stops: an obstruction}

Can the harmonic identity itself --- Theorem~\ref{thm:engine} with its $1/k$ weights --- be proved
by a plain sign-reversing involution on factorization tuples? No:

\begin{remark}[Obstruction]\label{rem:obstruction}
$\kappa$ is not integer-valued ($\kappa(p^2) = \tfrac12$), while any identity established by an
unweighted sign-reversing involution between sets of tuples computes an integer. A combinatorial
proof of Theorem~\ref{thm:engine} must therefore either pass through the M\"obius twin (Section~4)
or count something finer than tuples.
\end{remark}

The finer count has a natural candidate. The cyclic group $C_k$ acts on the $k$-tuples counted by
$D_k(m)$ by rotation, and orbit-counting rewrites the engine as
\[
\kappa(m) \;=\; \sum_{k \ge 1} (-1)^{k-1} \sum_{\text{orbits } O \text{ of } C_k}
\frac{1}{|\mathrm{Stab}(O)|},
\]
since $|O|/k = 1/|\mathrm{Stab}|$. The harmonic weights are now orbit-size corrections, nonunit
only on \emph{periodic} tuples. Micro-examples, machine-checked in Section~7: at $m = p^2$ the
$k = 2$ tuple $(p, p)$ is fixed by rotation, contributing $-\tfrac12$, and
$\kappa = 1 - \tfrac12 = \tfrac12$; at $m = pq$ the two tuples $(p, q), (q, p)$ form one free
orbit contributing $-1$, and $\kappa = 1 - 1 = 0$. This is the commutative shadow of the
Witt/necklace calculus of Metropolis and Rota \cite{MetropolisRota}, where logarithms count
aperiodic necklaces.

\begin{question}[The necklace question]\label{q:necklace}
Is there a weight-preserving sign-reversing involution on the set of pairs (cyclic orbit, marked
period) --- or another finer structure --- that proves Theorem~\ref{thm:engine} directly, with the
$1/j$ at axis powers emerging as the surviving periodic orbit? A positive answer would categorify
the engine; Section~6 records where such a categorification would live.
\end{question}

We record the honest outcome required by the work package: the bijective proof exists and is given
for the twin (Section~4); for the harmonic form the obstruction is Remark~\ref{rem:obstruction},
the reformulation is the display above, and the door is Question~\ref{q:necklace}.

\section{Categorification pointer}

At crystal generality the layers are already sets --- $D_k(m) = |L_k(m)|$ for the set $L_k(m)$ of
tuples --- so the layer identities categorify trivially. The substantive known categorification at
$C_\Z$ is the \emph{arithmetic product of species} of Maia and M\'endez \cite{MaiaMendez}, under
which Dirichlet series of species multiply; the tower of layers should be the tower of arithmetic
powers of the species of nonempty blocks. What would be genuinely new is a species-level
\emph{logarithm} producing $\kappa$ --- the categorified Witt formula of
Question~\ref{q:necklace}. We name the door and do not enter it in this note.

\section{Machine verification}

The script \texttt{formal\_engine.py} in the repository \cite{Code} checks every identity above by
exact rational arithmetic on truncated crystals, with no floating point and no analytic object
anywhere. Elements are exponent vectors; convolution is vector splitting; truncation is by
dimension $\Omega \le B$, which is exact for every identity here since all are $\Omega$-graded.
Five checks: (i) Theorem~\ref{thm:engine} on every element of random crystals with $r$ axes,
$\Omega \le B$, for $(r, B) = (4, 9), (6, 6), (1, 40)$; (ii) $\mu * \zeta = \delta$ and the
product formula for $\mu$; (iii) Proposition~\ref{prop:twin} with $\ell$ taken \emph{basis-valued}
--- $\ell(m)$ is the exponent vector itself in $\Q^r$, so one run verifies $\Lambda = \ell * \mu$
simultaneously for every Beurling length assignment on $r$ axes; (iv) $\exp(\kappa) = \zeta$ at
truncation; (v) at $C_\Z$, $\kappa(n)$ against the factorization oracle for all $n \le 3000$,
computed by integer Dirichlet convolutions. All pass; the run takes seconds and is wired into the
repository's \texttt{run\_all.py}.

\section{Revision plan for the companion paper}

The changes this note induces in \cite{CompanionI} are surgical. (1) After its Lemma~1, add a
remark citing this note: the identity is formal, holds at crystal generality, and needs no
convergence. (2) In its Section~3, the exactness assertion (denominator~$1$) gains the formal
explanation of Corollary~\ref{cor:pairing}. (3) In related work, add the incidence-algebra
attribution (Rota; Doubilet--Rota--Stanley) alongside Linnik, and Knopfmacher~\cite{Knopfmacher}
for the abstract setting. (4) No numerical content changes; no table changes. These edits are queued for the
Companion~I revision that accompanies the essay's v2 (milestone K1).

\section*{Disclosure of AI use}

This note was developed in an extensive collaboration with Claude (Anthropic; Claude Fable 5,
accessed through Claude Code, 2026). The direction of the work and all decisions about scope and
claims are the author's; the AI system drafted the text and proofs, wrote the verification script,
and assembled the attributions, the reason for its use being to enable a single non-specialist
author to develop, check, and document the material to a publishable standard. Every identity
claimed is machine-verified as described in Section~7. The author has reviewed the note and takes
full responsibility for its content.

\begin{thebibliography}{99}
\bibitem{CompanionI} C. Gribble, Prime counts and prime sums from multiplication tables: a
self-similar formula system, preprint, Zenodo (2026), \texttt{doi:10.5281/zenodo.21769103}.
\bibitem{CompanionII} C. Gribble, Certified exact sums of primes over dyadic intervals via a
ladder of generalized factorials, preprint, Zenodo (2026), \texttt{doi:10.5281/zenodo.21769105}.
\bibitem{CompanionIII} C. Gribble, Recovering the primes in a dyadic interval from exact moments:
two constructive routes and an obstruction theorem, preprint, Zenodo (2026),
\texttt{doi:10.5281/zenodo.21769107}.
\bibitem{Essay} C. Gribble, The multiplication crystal: dimension, diffraction, and the symmetries
of arithmetic, manuscript (2026), in the code repository \cite{Code}.
\bibitem{Defs} C. Gribble, Crystals and their shadows: definitions for a programme on the
multiplication crystal, programme note (2026), in the code repository \cite{Code}.
\bibitem{Linnik} Yu.~V. Linnik, \emph{The Dispersion Method in Binary Additive Problems}, Amer.
Math. Soc., Providence, 1963.
\bibitem{Rota} G.-C. Rota, On the foundations of combinatorial theory I: Theory of M\"obius
functions, \emph{Z. Wahrscheinlichkeitstheorie} \textbf{2} (1964), 340--368.
\bibitem{DRS} P. Doubilet, G.-C. Rota, and R. Stanley, On the foundations of combinatorial theory
VI: The idea of generating function, in \emph{Proc. Sixth Berkeley Symp. Math. Statist. Prob.},
vol.~II, Univ. California Press, 1972, 267--318.
\bibitem{StanleyEC} R.~P. Stanley, \emph{Enumerative Combinatorics}, vol.~1, 2nd ed., Cambridge
Univ. Press, 2012.
\bibitem{MetropolisRota} N. Metropolis and G.-C. Rota, Witt vectors and the algebra of necklaces,
\emph{Adv. Math.} \textbf{50} (1983), 95--125.
\bibitem{MaiaMendez} M. Maia and M. M\'endez, On the arithmetic product of combinatorial species,
\emph{Discrete Math.} \textbf{308} (2008), 5407--5427.
\bibitem{Knopfmacher} J. Knopfmacher, \emph{Abstract Analytic Number Theory}, North-Holland,
Amsterdam, 1975.
\bibitem{Code} C. Gribble, Reference implementations accompanying this programme, code repository
(2026), \texttt{https://github.com/carlgribble-caa/prime-moments}.
\end{thebibliography}

\end{document}
